\subsection{Exploratory follow-up: realized descriptors and bounded regression}
\label{sec:bounded_method}
An additional historical comparison distinguishes the original generator-based
descriptor values, $X^{\rm G}$, from 22 descriptors recomputed on the loaded
instance and constructed MILP, $X^{\rm M}$. The latter use input profiles,
capacities, terminal-storage constraints and declared model dimensions before
presolve. They use no incumbent, optimized flow, solver-search statistic or
outcome as an explanatory variable. Coordinate names A1--B10 retain their
positions for correspondence, but several definitions change; this is a
separate representation rather than a silent correction of the original index.
In particular, A3 uses the actual aggregate demand peak and valley, B3 counts
unfixed binary variables before presolve, B8 uses realized storage durations
and terminal windows, and $B9^{\rm M}=|\mathcal L|/N$ uses actual lines.
The original raw values precede the reference normalization, while retaining
any clipping intrinsic to their original formulas.

For each representation, feature means and population standard deviations
are learned only on the relevant training fold, including every inner fold:
\begin{equation}
 \widetilde x_{sj}=(x_{sj}-\mu_j^{\rm train})/\sigma_j^{\rm train}.
\end{equation}
A training-constant coordinate receives a zero coefficient, even if it varies
at test time. Both representations use the same historical observations,
preferences, five deadlines and scenario folds. The projected comparator in
this follow-up is the saved deadline-specific positive linear predictor on
these standardized raw features,
\begin{equation}
 \widehat Y^{\rm proj}_{sb}=\min\{1,\max\{0,a_b+\boldsymbol\beta_b^\top
 \widetilde{\boldsymbol x}_s\}\},\qquad\boldsymbol\beta_b\geq0.
 \label{eq:followup_projection}
\end{equation}
It is distinct from the reference-normalized models in
Table~\ref{tab:revision_models}; their historical losses are preserved.
Projection cannot increase squared error when the response is in $[0,1]$.

After inspecting the realized-descriptor extrapolations on the historical
cohort, we fitted three explicitly exploratory alternatives: positive linear
regression with B9 removed and predictions projected as above, positive
fractional-logistic regression with all coordinates, and the same bounded
regression with B9 removed. Removal deletes only $B9^{\rm G}$ or $B9^{\rm M}$
from the respective feature set; all remaining coefficients and the penalty
are refitted. For preference $p$ and deadline $b$, the bounded model is
\begin{align}
 \eta_s&=a_b^{(p)}+\boldsymbol\beta_b^{(p)\top}\widetilde{\boldsymbol x}_s,
 &\widehat Y^{(p)}_{sb}&=(1+\exp(-\eta_s))^{-1},\nonumber\\
 (\widehat a,\widehat{\boldsymbol\beta})
 &=\underset{a,\,\boldsymbol\beta\geq0}{\operatorname{argmin}}\,
 \left\{\frac1n\sum_{s=1}^n
 [\log(1+\exp(\eta_s))-Y^{(p)}_{sb}\eta_s]
 +\frac{\kappa}{2}\|\boldsymbol\beta\|_2^2\right\}.
 \label{eq:bounded_fractional}
\end{align}
The continuous responses are not rounded to binary labels. The output
estimates a mean operational loss, not a binary failure probability.
The intercept is unpenalized and constant-response training folds use their
observed constant. This changes the link and fitting loss jointly, retaining
the positivity assumption; it is not an isolated test of the link function.

The same candidate penalties $\{0,10^{-4},10^{-3},10^{-2},10^{-1},1,10\}$
are selected inside the training data. Five outer scenario folds use three
inner scenario folds. Six family-held-out evaluations use the five remaining
families as inner folds; no family indicator is a predictor. The linear
ablation retains the original unprojected inner-MSE selection criterion,
whereas fractional-logistic tuning uses the MSE of its bounded predictions.
Final evaluation uses delivered bounded predictions for every arm. Thus the
linear ablation matches the original fitting procedure, while the bounded
model is a different predictive procedure. The design comprises 330 saved
outer-fold fits, each containing five deadline models, or 1,650 models in total.

Paired loss contrasts average the five deadlines within each scenario before
2,000 scenario bootstrap resamples. Within each validation scheme, 40
simultaneous contrasts cover three alternatives against the projected
comparator and bounded B9 removal against the full bounded model, for both
representations and all five preferences. These conditional intervals do not
include refitting uncertainty or adjust for choosing this follow-up after
examining historical results. Neither nested folds nor simultaneous intervals
make this a confirmatory independent test.

This follow-up changes no physical dispatch, historical reference, E/ES
classifier or time-allocation policy. Constructing the realized descriptors
takes a median 1.02 seconds per historical scenario and is outside the
archived operational deadlines. Predictive gains therefore do not establish
net online decision gains.
